\chapter{Algoritmo genetico OpenAI}

Il tool progettato permette di eseguire l'ottimizzazione della funzione obbiettivo $\mathcal{L}$
tramite diversi ottimizzatori come Nevergrad e Nomad. Può essere facilmente esteso per includere diversi
algoritmi.

In questa relazione di tirocinio si propone l'implementazione di un algoritmo genetico di OpenAI \cite{salimans2017es} per risolvere problemi
blackbox.

Un algoritmo genetico è una procedura di ricerca euristica ispirate all'evoluzione naturale osservabile in natura.

Ad ogni iterazione dell'algoritmo (generazione) un insieme di parametri (genotipi) che vengono perturbati (mutazione genetica) e la funzione obbiettivo
viene valutata su queste perturbazioni (fitness).

Le perturbazioni con la fitness migliori vengono ricombinate per formare i nuovi parametri per la nuova generazione.

In questo algoritmo di OpenAI la popolazione è rappresentata come una distribuzione gaussiana multivariata che ha media sui parametri.

Quindi si voglio trovare i parametri che  massimizzano:

\begin{equation}
    \mathbb{E}_{\varepsilon \sim \mathcal{N}(0,I)} \{ \mathcal{L}(\theta + \sigma \varepsilon) \}
\end{equation}

Possiamo ottimizzare i parametri $\theta$ usando un gradient descent stocastico, vedere la figura \ref{gradient_descend} per una visualizzazione grafica del funzionamento.

\begin{equation}
    \nabla_{\theta} \mathbb{E}_{\varepsilon \sim \mathcal{N}(0,I)} \{ \mathcal{L}(\theta + \sigma \varepsilon) \} = -\frac{1}{\sigma} \mathbb{E}_{\varepsilon \sim \mathcal{N}(0,I)} \{ \mathcal{L}(\theta + \sigma \varepsilon) \varepsilon \}
\end{equation}


Esso può essere stimato campionando $n$ valori $\varepsilon_1, \dots, \varepsilon_n \sim \mathcal{N}(0,I)$.

\begin{equation}
    -\frac{1}{\sigma} \mathbb{E}_{\varepsilon \sim \mathcal{N}(0,I)} \{ \mathcal{L}(\theta + \sigma \varepsilon) \varepsilon \} \approx -\frac{1}{\sigma n} \sum_{i=1}^n \mathcal{L}(\theta + \sigma \varepsilon_i) \varepsilon_i
\end{equation}

Ad ogni generazione i nuovi parametri si spostano in direzione del gradiente.

Poiché il gradiente può avere grande ampiezza ,esso va scalato per un iperparametro dell'algoritmo $\alpha > 0$ detto \textbf{learning rate}.

\begin{figure}
\centering
    \begin{tikzpicture}
    \begin{axis}[width=12cm,%
        declare function={f(\x,\y)=cos(deg(\x)*0.8)*cos(deg(\y)*0.6)*exp(0.1*\x);}]
    \addplot3[surf,shader=interp,domain=-4:4,%samples=81
    ]{f(x,y)};
    \edef\myx{0.15} % first x coordinate
    \edef\myy{-0.15} % first y coordinate
    \edef\mystep{-2}% negative values mean descending
    \pgfmathsetmacro{\myf}{f(\myx,\myy)}
    \edef\lstCoords{(\myx,\myy,\myf)}
    \pgfplotsforeachungrouped\X in{0,...,5}
    {
    \pgfmathsetmacro{\myx}{\myx+\mystep*xgrad(\myx,\myy)}
    \pgfmathsetmacro{\myy}{\myy+\mystep*ygrad(\myx,\myy)}
    \pgfmathsetmacro{\myf}{f(\myx,\myy)}
    \edef\lstCoords{\lstCoords\space (\myx,\myy,\myf)}
    }
    \addplot3[samples y=0,arrowed] coordinates \lstCoords;
    \end{axis}
    \end{tikzpicture}
\caption[Visualizzazione del gradient descent, presa da qui]{Visualizzazione del gradient descent, presa da \href{https://tex.stackexchange.com/questions/544796/plot-gradient-descent}{https://tex.stackexchange.com/questions/544796/plot-gradient-descent}}
\label{gradient_descend}
\end{figure}

    
\begin{algorithm}
\caption{OpenAI Algoritmo Evolutivo Genetico}\label{alg:openai_algo}
\begin{algorithmic}[1]
\State \textbf{Input: } Learning rate $\alpha$, deviazione standard $\sigma$, parametri iniziali: $\theta_0$
\For{$t = 0, 1, 2, \dots$}
    \State Campiona $\varepsilon_1, \dots, \varepsilon_n \sim \mathcal{N}(0,I)$
    \State Calcola funzione di loss $\mathcal{L}_i = \mathcal{L}(\theta_t + \sigma \varepsilon_i)$ \textbf{for} $i \in [1,n]$
    \State $\theta_{t+1} \gets \theta_t - \alpha \frac{1}{\sigma n} \sum_{i=1}^n  \mathcal{L}_i \varepsilon_i$
\EndFor
\end{algorithmic}
\end{algorithm}
